% ===== 原文 pp.41–42 =====
% 译者 tr-06：§5.1 后半（谱曲线、定理 5.5）、注 5.6、§5.2 开头（定义 5.7、例 5.8、注 5.9）
% 编号区间：式 (5.4)、定理 5.5、注 5.6、定义 5.7、例 5.8、注 5.9
% 依据：split/p041.txt、split/p042.txt（仅用于定位）＋ pages/p041.png、p042.png（200 dpi，数学逐符号核对）

% [原文 p.41]
为了理解经由群同态得到的 Higgs 丛之间特征值的关系，回顾如下事实是有用的：Higgs 场 $\Phi$ 的特征多项式在典范丛 $K$ 的全空间中定义一条曲线，其方程可以通过考虑
\[
p(x)=x^n+\tr(\Phi)x^{n-1}+\sum_{k=2}^{n-1}x^{n-k}(-1)^k\Tr(\Lambda^k\Phi)+(-1)^n\det(\Phi)
\]
而得到；其中 $\Tr(\Lambda^k\Phi)$ 是 $\Phi$ 的 $k$ 次外幂的迹，其维数为 $\binom{n}{k}$，与前面的式 (3.2) 相同。

% [Tutor Note: p.41 此式（与 p.15 的式 (3.2) 完全相同）中 $\tr(\Phi)x^{n-1}$ 前的符号应为负号：按同一公式自身的通项模式，$\det(I\eta-\Phi)=\sum_{k=0}^{n}(-1)^k\Tr(\Lambda^k\Phi)\,\eta^{n-k}$，其 $k=1$ 项即 $-\tr(\Phi)\eta^{n-1}$，这同时也是 (5.4) 的展开。实用上不受影响，因为对 $\SL(n,\C)$ 有 $\tr(\Phi)=0$，故 (5.4) 的级数直接从 $a_2\eta^{n-2}$ 起。译文照原样排正号；同一式中 $\tr$ 与 $\Tr$ 混用亦照原样。]

如此前所见，并如 [61] 所解释的，与 $\SL(n,\C)$-Higgs 丛 $(E,\Phi)$ 相伴的谱曲线 $\pi : S\to\Sigma$ 的方程为
\begin{equation}
S := \{0=\det(I\eta-\Phi)\}=\{0=\eta^n+a_2\eta^{n-2}+\dots+a_{n-1}\eta+a_n\}, \tag{5.4}
\end{equation}
其中 $a_i\in H^0(\Sigma,K^i)$，而 $\eta$ 是 $\pi^*K$ 的 重言截面。谱曲线 $S$ 一般光滑，Hitchin 底中一般点处的纤维可等同于 Prym 簇 $\Prym(S,\Sigma)$。固定 $K^{1/2}$ 的一个选择后，向量丛 $E$ 可由 $L\in\Prym(S,\Sigma)$ 还原为 $E := \pi_*(L\otimes\pi^*(K^{(n-1)/2}))$。

对 $\SO(2n,\C)$-Higgs 丛，Higgs 场 $\Phi$ 定义覆盖 $\pi : S\to\Sigma$，其方程由式 (5.4) 给出，其中 $a_n = p_n\in H^0(\Sigma,K^n)$ 是 $\Phi$ 的 Pfaffian。该曲线在 $\eta = p_n = 0$ 处奇异，这些点是一个自然对合 $\sigma : \eta\mapsto-\eta$ 的不动点。$S$ 的正规化 $\hat\pi : \hat S\to\Sigma$ 就是我们所称的谱曲线。这条曲线 $\hat S$ 一般光滑，且对合 $\sigma$ 延拓为无不动点的对合 $\hat\sigma$。一般的 $\SO(2n,\C)$ 纤维由 $\Prym(\hat S,\hat S/\sigma)$ 的一个连通分支给出。与前一情形类似，选定 $K^{1/2}$ 后，向量丛 $E$ 可由 $L\in\Prym(\hat S,\hat S/\sigma)$ 还原为
\[
E := \pi_*(L\otimes(K_{\hat S}\otimes\pi^*(K^*))^{1/2}).
\]

% [Tutor Note: p.41 上一句刚把延拓后的无不动点对合记作 $\hat\sigma$，故紧接着两处 $\Prym(\hat S,\hat S/\sigma)$ 的分母应为 $\hat S/\hat\sigma$；原文印作 $\sigma$，译文照原样。可核对：$S$ 在 $\eta=p_n=0$ 处的奇异点正是 $\sigma:\eta\mapsto-\eta$ 的不动点，正规化后 $\hat\sigma$ 无不动点，于是 $\hat S/\hat\sigma$ 光滑，符合 Prym 簇的定义要求（$\hat S\to\hat S/\hat\sigma$ 为无分歧二重覆盖）。]

借助上述谱数据，连同命题 5.3 所给出的对应，同源诱导出 [29, 定理 1 与定理 2] 中 Hitchin 纤维化上的如下映射。此外，定理 5.5 中的态射可限制到实的情形（例如见问题 16 与 17）。

\begin{thm}{5.5（[29]）}
考虑两个一般的 $\SL(2,\C)$-Higgs 丛 $(E_i,\Phi_i)$，其谱数据为 $S_i$ 与 $L_i\in\Prym(S_i,\Sigma)$。则 $\cI_2$ 是到其像上的 $2^{2g}$ 重覆盖，而像 $\cI_2((E_1,\Phi_1),(E_2,\Phi_2))$ 的谱数据由 $(\hat S_4,\cL_4)$ 给出，其中
\begin{itemize}[leftmargin=1.4em,itemsep=1pt,topsep=1pt]
\item $\hat S_4 := S_1\times_{\Sigma}S_2$ 是光滑的分歧四重覆盖，且
\item $\cL_4 := p_1^*(L_1)\otimes p_2^*(L_2)$，其中 $p_i : S_1\times_{\Sigma}S_2\to S_i$ 是投影映射。
\end{itemize}
此外，对一个一般的 $\SL(4,\C)$-Higgs 丛 $(E',\Phi')$（谱数据为 $(S',L')$），由像 $\cI_3(E,\Phi)$ 给出的 $\SO(6,\C)$-Higgs 丛的谱数据为 $(\hat S_6,\cI_3(L))$，其中
\begin{itemize}[leftmargin=1.4em,itemsep=1pt,topsep=1pt]
\item $\hat S_6$ 是 $S'\times_{\Sigma}S'$ 中非对角分支的对称化；
\item $\cI_3(L)$ 是由 $\cL' := p_1^*(L)\otimes p_2^*(L)$ 的局部截面生成的典范扭（canonical twist），这些局部截面关于 $S'\times_{\Sigma}S'$ 的对称是反不变的。
\end{itemize}
\end{thm}

% [Tutor Note: p.41 定理 5.5 中谱数据的线丛记号不统一：$L_i$、$L_1$、$L_2$、$L'$ 用正体 $L$，而 $(\hat S_4,\cL_4)$、$\cL_4:=\dots$、$\cL'$ 用花体 $\cL$。译文按原样保留两种字形未作归一（若全书统一，宜另行裁定后统一替换）。同一段还有一处指代不一致：原文说 "for a generic $\SL(4,\C)$-Higgs bundle $(E',\Phi')$ ... the image $\cI_3(E,\Phi)$ ... $\cL' := p_1^*(L)\otimes p_2^*(L)$"，其中的 $E,\Phi,L$ 按上下文应为 $E',\Phi',L'$。译文照原样。]

% [原文 p.42]
应当注意，定理 5.5 所考虑的 Higgs 丛全都落在 Hitchin 纤维化的一般纤维之内。然而，该同源可下降到这些复 Lie 群的所有实形式，因此人们应当能够得到完全位于 Hitchin 纤维化奇异轨迹之上的实 Higgs 丛之间的态射。最早被研究的情形涉及 ribbon 的几何，见 [97]。

\begin{rem}{5.6}
[29] 的方法在研究 F-理论中与椭圆纤维化 CY 三维簇（elliptically fibred CY 3-folds）的退化有关的问题时也同样有用，[4] 即如此做了。
\end{rem}

\subsection{多边形与超多边形}

当考虑抛物 Higgs 丛（如第 2.3 节所定义）与箭图簇（quiver varieties，即箭图的表示簇，此时底流形是一个点）时，数学对象之间会出现一种有趣的关系；[77,78] 以及许多研究者随后的工作对此作了研究。

\begin{defn}{5.7}
一个箭图（quiver）$Q=(V,A,h,t)$ 是一个有向图，由有限的顶点集 $V$、有限的箭向集 $A$，以及头、尾映射 $h,t : A\to V$ 组成。
\end{defn}

\begin{exmp}{5.8}
箭图可以用有向图来表示；在后续各节中将会显得有用的一个例子是如下星形箭图（star shaped quiver），其中顶点集为 $V=\{v_1,\dots,v_3\}$，箭向集为 $A=\{a_i : v_i\to v_0\mid 1\leqslant i\leqslant 3\}$：
\begin{center}
\begin{tikzpicture}[x=0.8cm,y=0.8cm]
\fill (0,1.5) circle (1.6pt);
\node at (0.34,1.5) {$v_1$};
\fill (0,0) circle (1.6pt);
\node at (0.34,0) {$v_0$};
\fill (-2.6,0) circle (1.6pt);
\node at (-2.94,0) {$v_3$};
\fill (2.6,0) circle (1.6pt);
\node at (2.94,0) {$v_2$};
\draw[->] (0,1.34) -- (0,0.18);
\node at (-0.48,0.76) {$a_1$};
\draw[->] (-2.4,0) -- (-0.18,0);
\node at (-1.3,-0.42) {$a_3$};
\draw[->] (2.4,0) -- (0.18,0);
\node at (1.3,-0.42) {$a_2$};
\end{tikzpicture}
\end{center}
这里，头映射由 $h(a_i)=v_0$ 给出，尾映射由 $t(a_i)=v_i$ 给出。
\end{exmp}

% [Tutor Note: p.42 原文 "as bellow" 应为 "as below"（原文笔误）。]
% [Tutor Note: p.42 例 5.8 中的星形箭图是原文唯一的未编号行内箭图（并非 Figure 5.10，不属于两张编号图）。为不丢内容，译文用最小化的 TikZ 重绘：中心 $v_0$；外围 $v_1$（上方）、$v_3$（左）、$v_2$（右）；$a_1$ 自上而下指向 $v_0$，$a_3$ 自 $v_3$ 向右指向 $v_0$，$a_2$ 自 $v_2$ 向左指向 $v_0$；不设图号与图注。若编辑部希望改用原图裁剪（需新增 figures 条目），可整段替换。]

箭图 $Q$ 的维数向量（dimension vector）是正整数构成的一个 $(n+1)$-元组 $d=(r,1,\dots,1)$，这个向量把 $Q$ 唯一地确定为星形箭图。此外，箭图自然地带有定义好的表示空间：
\[
\Rep(Q)=\bigoplus_{i=1}^{n}\Hom(\C,\C^r)\cong\Hom(\C^n,\C^r)=\C^{n\times r}.
\]

% [Tutor Note: p.42 这里的维数核对：$\Hom(\C,\C^r)\cong\C^r$，$n$ 个直和给出 $\C^{nr}=\C^{n\times r}$，与 $\Hom(\C^n,\C^r)$ 一致；而 $\dim\Rep(Q)=nr$ 正对应维数向量 $d=(r,1,\dots,1)$（中心顶点维数 $r$，$n$ 个外围顶点各维数 $1$）。]

我们特别感兴趣的是星形箭图 $Q_n$：它有一个中心顶点 $v_0$ 与 $n$ 个外围顶点 $v_0,\dots,v_n$，以及箭向 $a_i : v_i\to v_0$，如例 5.8 中 $n=3$ 的那一个。下文我们将考虑固定的维数向量 $d=(r,1,\cdots,1)\in\N^{n+1}$，并记住这一取值，把该星形箭图记作 $Q^r_n$。

% [Tutor Note: p.42 原文 "one central vertex $v_0$ and $n$ outer vertices $v_0,\dots,v_n$" 中第一个 $v_0$ 应为 $v_1$：中心顶点既已记作 $v_0$，$n$ 个外围顶点应为 $v_1,\dots,v_n$。依据：同句后文 "arrows $a_i : v_i\to v_0$"、例 5.8 的 $V=\{v_1,\dots,v_3\}$（$n=3$），以及 $d=(r,1,\dots,1)$ 是 $(n+1)$ 元组（$r$ 对应中心 $v_0$，$n$ 个 $1$ 各对应一个 $v_i$，$i\geqslant1$）。译文照原样排。]

\begin{rem}{5.9}
箭图 $Q$ 的一个表示也可以写成
\[
W := ((W_v)_{v\in V},(\varphi_a)_{a\in A}),
\]
其中 $W_v$ 是有限维向量空间；且对所有的 $a\in A$，$\varphi_a : W_{t(a)}\to W_{h(a)}$ 是线性映射。
\end{rem}
